\documentclass[10pt,letterpaper]{article}
\usepackage{cornell-notes}

\title{Chapter 10: Random Integer Partition}
\author{}
\date{}

\setDocTitle{Chapter 10: Random Integer Partition}
\setDocSubtitle{Combinatorial Algorithms}
\setDocVersion{v1.0}
\setDocStatus{Study Companion}
\setDocOwner{Combinatorial Algorithms Cornell Notes}
\setDocAuthor{}
\setDocDate{\today}
\setDocSummary{Cornell-format study and review notes for Chapter 10: Random Integer Partition.}

\setCornellCollection{Combinatorial Algorithms Cornell Notes}
\setCornellUnitType{Chapter}
\setCornellUnitNumber{10}
\setCornellUnitTitle{Random Integer Partition}
\setCornellCompanionLabel{Study and review companion}

\hypersetup{pdftitle={Chapter 10: Random Integer Partition},pdfsubject={Cornell notes},pdfauthor={}}

\begin{document}
\maketitle
\begin{tcolorbox}[colback=Paper,colframe=Ink]{\small\bfseries\color{Teal}CORNELL NOTES}\par{\LARGE\bfseries Chapter 10: Random Partition of an Integer $n$ (RANPAR)}\par\medskip\textbf{Topic:} Recurrence-driven uniform sampling\hfill\textbf{Date:} \today\par\textbf{Study purpose:} Turn Euler's partition recurrence into sampling probabilities.\end{tcolorbox}
\begin{tcolorbox}[colback=Teal!5,colframe=Teal,title=Essential question]How can a counting identity be normalized into transition probabilities that construct every partition of $n$ with probability $1/p(n)$?\end{tcolorbox}
\begin{tcolorbox}[colback=white,colframe=Mist,title=Learning objectives]\begin{itemize}\item Explain the largest-part recurrence alternative.\item State Euler's divisor-sum identity.\item Normalize the identity into probabilities for pairs $(d,j)$.\item Trace the residual-size process.\item Summarize the uniformity proof and required precomputation.\end{itemize}\textbf{Prerequisites:} integer partitions, divisor sums, recurrence relations, conditional probability.\end{tcolorbox}
\section*{Cornell notes}
\begin{longtable}{@{}Q!{\color{Mist}\vrule width 1pt}A@{}}\cornellhead
\cue{What is the target distribution?}{Choose one of the $p(n)$ integer partitions with probability $1/p(n)$. The output is stored by multiplicity: $\operatorname{MULT}(d)$ is the number of copies of part $d$.}
\cue{What first recurrence is discussed?}{If $p(n,k)$ counts partitions of $n$ with largest part $k$, then
\[
p(n,k)=p(n-1,k-1)+p(n-k,k).
\]
This supports a recursive sampler but requires a two-index table. The implemented method instead uses a one-dimensional table of $p(m)$.}
\cue{What identity drives RANPAR?}{With $\sigma(r)$ the sum of positive divisors of $r$,
\[
n,p(n)=\sum_{m=0}^{n-1}\sigma(n-m)p(m)
=\sum_{d\ge1}\sum_{j\ge1}d\,p(n-jd).
\]
Terms with negative arguments are zero. The identity is proved both from the partition generating function and by a multiple-counting construction.}
\cue{How are probabilities obtained?}{At residual total $m$, choose a pair $(d,j)$ with $jd\le m$ and
\[
\Pr(d,j\mid m)=\frac{d\,p(m-jd)}{m,p(m)}.
\]
The normalized identity guarantees these nonnegative probabilities sum to one.}
\cue{What is the transition?}{Add $j$ copies of part $d$ to the output, then set $m\leftarrow m-jd$. Repeat until $m=0$. Every step strictly reduces the residual, so termination is immediate.}
\cue{Why is the final partition uniform?}{Fix a partition with distinct parts $d_i$ and multiplicities $\mu_i$. It can be reached by choosing any $j\in\{1,\ldots,\mu_i\}$ copies of any $d_i$ at the first step. By induction, the remaining partition has probability $1/p(n-jd_i)$. Substitution cancels the recurrence factors, and summing $d_i$ over all its $\mu_i$ copies gives total probability $1/p(n)$.}
\cue{What is precomputed?}{The routine stores $p(1),\ldots,p(n)$ and extends the table only when a larger $n$ is requested. Euler's identity supplies each next value. The same table then weights pair selection at every residual size.}
\cue{State and cost cautions}{Partition counts grow rapidly; the source uses integer storage appropriate to its tested range, but a modern implementation must guard overflow or use arbitrary precision. Random thresholds should be formed without losing probability mass. Multiplicity totals must satisfy $\sum_d d\,\operatorname{MULT}(d)=n$.}
\cue{Historical or legacy context}{The routine keeps its largest precomputed argument in static procedure state and uses a global-style random function. Modern code should make the count table and random generator explicit dependencies.}
\cue{Retrieval practice}{\begin{enumerate}\item Why prefer the one-index identity?\item What is $\sigma(r)$?\item What is the probability of $(d,j)$?\item Why must the residual terminate?\item How does the induction cancel $p(m-jd)$?\end{enumerate}}
\end{longtable}
\begin{tcolorbox}[colback=Ink!4,colframe=Ink,title=Cornell summary]
RANPAR illustrates a general sampling method: prove a positive recurrence for counting objects, divide by the total count, and interpret the summands as transition probabilities. Euler's identity expresses $m p(m)$ as a sum of $d p(m-jd)$. At residual $m$, the algorithm samples $(d,j)$ with the normalized weight, appends $j$ copies of $d$, and reduces the residual by $jd$. Induction over the residual and a sum over the possible first groups of equal parts prove that every final partition has probability $1/p(n)$. A one-dimensional table of partition counts supports both precomputation and sampling, but count growth and numeric overflow require explicit attention in modern implementations.
\end{tcolorbox}
\end{document}
